\subsection{Local Hermite regression scheme}

Let
\[
\pi=\{0=t_0<\cdots<t_N=T<t_{N+1}<\cdots<t_{N'}=T+\delta\}
\]
be a uniform grid with time step $h=t_{n+1}-t_n$, chosen so that
$\delta/h$ is an integer. We simulate $M$ independent Brownian paths
$(W_{t_n})_{0\leq n\leq N'}$. On the terminal interval $[T,T+\delta]$,
the prescribed processes are known and we set
\[
Y^\pi_{t_n}=\Phi(t_n,W_{t_n}),
\qquad
Z^\pi_{t_n}=\Psi(t_n,W_{t_n}),
\qquad N\leq n\leq N'.
\]
These terminal data are the only oracle values used by the backward algorithm.
In what follows, the superscript $h$ denotes the theoretical time-discrete
scheme, while its grid values used in the regressions are denoted by the
superscript $\pi$.

We first recall the time-discrete scheme of Section~\ref{se:approx}. Let
$\overline Y^h$ and $\overline Z^h$ denote the piecewise-constant extensions of
the grid values. At each date $t_n<T$, the scheme reads
\[
\widetilde Y^h_{t_n}
=
\mathbb E_{t_n}\!\left[Y^h_{t_{n+1}}\right],
\qquad
Z^h_{t_n}
=
\frac{1}{h}\mathbb E_{t_n}\!\left[
Y^h_{t_{n+1}}\Delta W_n
\right],
\]
and
\[
Y^h_{t_n}
=
\widetilde Y^h_{t_n}
+h f\!\left(
t_n,X^h_{t_n},Y^h_{t_n},Z^h_{t_n},\mu^h_{t_n}
\right),
\qquad
\mu^h_{t_n}=\mu^{\overline Y^h,\overline Z^h}_{t_n},
\]
where
\[
\Delta W_n=W_{t_{n+1}}-W_{t_n}.
\]
The implementation simplifies whenever the anticipated term $\mu^h_{t_n}$ does
not depend on the current value $Y^h_{t_n}$. Indeed, future grid values are
already available in the backward recursion and $Z^h_{t_n}$ is computed before
the anticipated term. If, in addition, the driver also does not depend on
$Y^h_{t_n}$, the last update is explicit and no fixed-point iteration is
required.

The three examples below have precisely this sequential structure. Moreover,
their drivers depend on the measure $\mu^h_{t_n}$ only through a conditional
expectation, which we denote by
$A^h_{t_n}:=\mathbb{E}_{t_n}[\mathcal{A}_n]$. For
$n=N-1,\ldots,0$, we therefore approximate sequentially the four quantities
\[
\widetilde Y^h_{t_n},
\qquad
Z^h_{t_n},
\qquad
A^h_{t_n},
\qquad
Y^h_{t_n}.
\]

All four quantities are approximated with the same degree-three Hermite model.
For $t_n>0$, set
\[
X_n=\frac{W_{t_n}}{\sqrt{t_n}}
\]
and define
\[
H_0(x)=1,
\qquad
H_1(x)=x,
\qquad
H_2(x)=\frac{x^2-1}{\sqrt 2},
\qquad
H_3(x)=\frac{x^3-3x}{\sqrt 6}.
\]
For $\beta=(\beta_0,\ldots,\beta_3)\in\mathbb R^4$, write
\[
P_\beta(x)
=
\sum_{k=0}^3\beta_k H_k(x).
\]
At $t_0=0$, the nonconstant basis functions are set to zero, so that the fitted
quantity is constant, as required by $\mathcal F_0$-measurability.

The four parameter vectors are determined sequentially,
\[
\widehat\beta_n^{\widetilde Y}
=
\arg\min_{\beta\in\mathbb R^4}
\mathbb E\!\left[
\left|P_\beta(X_n)-Y^\pi_{t_{n+1}}\right|^2
\right],
\qquad
\widehat{\widetilde Y}_n=P_{\widehat\beta_n^{\widetilde Y}}(X_n),
\]
\[
\widehat\beta_n^Z
=
\arg\min_{\beta\in\mathbb R^4}
\mathbb E\!\left[
\left|
P_\beta(X_n)
-
\frac{
\bigl(Y^\pi_{t_{n+1}}-\widehat{\widetilde Y}_n\bigr)\Delta W_n}{h}
\right|^2
\right],
\qquad
\widehat Z_n=P_{\widehat\beta_n^Z}(X_n),
\]
\[
\widehat\beta_n^A
=
\arg\min_{\beta\in\mathbb R^4}
\mathbb E\!\left[
\left|P_\beta(X_n)-\mathcal A_n\right|^2
\right],
\qquad
\widehat A_n=P_{\widehat\beta_n^A}(X_n),
\]
\[
\widehat\beta_n^Y
=
\arg\min_{\beta\in\mathbb R^4}
\mathbb E\!\left[
\left|
P_\beta(X_n)
-
\left(
\widehat{\widetilde Y}_n
+h f\!\left(t_n,P_{\widehat\beta_n^A}(X_n)\right)
\right)
\right|^2
\right],
\qquad
\widehat Y_n=P_{\widehat\beta_n^Y}(X_n).
\]

The centered target in the $Z$ criterion is equivalent to
$h^{-1}Y^\pi_{t_{n+1}}\Delta W_n$, since
$\widehat{\widetilde Y}_n$ is $\mathcal F_{t_n}$-measurable and
$\mathbb E_{t_n}[\Delta W_n]=0$. The centering reduces the variance of the
regression target.

All four problems are instances of the same ridge regression. For a
generic sample target $(R_n^{(m)})_{1\leq m\leq M}$, we compute
\[
\widehat\beta_n[R]
=
\arg\min_{\beta\in\mathbb R^4}
\left\{
\sum_{m=1}^M
\left|
P_\beta(X_n^{(m)})-R_n^{(m)}
\right|^2
+\lambda|\beta|^2
\right\},
\qquad \lambda=10^{-10}.
\]
The normal equations are solved in double precision. The small ridge parameter
is used only to stabilize the linear system.

\subsection{Benchmark examples}

We test the regression scheme on three explicit one-dimensional examples adapted
from \citet{peng2009anticipated}. In each case, the closed-form solution is used
only for out-of-sample evaluation.

\paragraph{Example 4.3 of \citet{peng2009anticipated}.}
For $0\leq t\leq T$,
\[
Y_t
=
T W_T
-
\int_t^T \frac{1}{s+\delta}\,
\mathbb E_s[Y_{s+\delta}]\,ds
-
\int_t^T Z_s\,dW_s,
\]
with $Y_t=tW_t$ on $[T,T+\delta]$. The exact solution is
\[
Y_t=tW_t,
\qquad
Z_t=t,
\qquad 0\leq t\leq T+\delta.
\]

\paragraph{Example 5.3 (I) of \citet{peng2009anticipated}.}
For $0\leq t\leq T$,
\[
Y_t
=
W_T^2-T
-
\int_t^T
\sqrt{\frac{\pi}{2\delta}}\,
\mathbb E_s\!\left[|Z_{s+\delta}-Z_s|\right]ds
-
\int_t^T Z_s\,dW_s,
\]
with
\[
Y_t=W_t^2-T-(T-t),
\qquad
Z_t=2W_t,
\qquad t\in[T,T+\delta].
\]
The exact solution, on the full interval, is
\[
Y_t=W_t^2-T-(T-t),
\qquad
Z_t=2W_t.
\]

\paragraph{Example 5.3 (II) of \citet{peng2009anticipated}.}
For $0\leq t\leq T$,
\[
Y_t
=
4W_T^2
-
\int_t^T
\sqrt{\frac{\pi}{2\delta}}\,
\mathbb E_s\!\left[|Z_{s+\delta}-Z_s|\right]ds
-
\int_t^T Z_s\,dW_s,
\]
with
\[
Y_t=4W_t^2-4(T-t),
\qquad
Z_t=8W_t,
\qquad t\in[T,T+\delta].
\]
The exact solution is
\[
Y_t=4W_t^2-4(T-t),
\qquad
Z_t=8W_t.
\]

Figures~\ref{fig:example43_quantiles}, \ref{fig:example531_quantiles}, and
\ref{fig:example532_quantiles} compare the $5\%$, $50\%$, and $95\%$ quantiles
of the learned processes with their exact counterparts. Solid lines denote exact
quantiles and hollow circles denote Hermite regression estimates. The markers lie
on the corresponding lines in all three examples. In particular, the method
recovers the deterministic process $Z_t=t$ in Example~4.3 even though the
regressions are functions of the Brownian state.

The quantile figures use a representative model trained with $N=20$ and
$10^6$ Brownian trajectories; the displayed quantiles are computed from an
independent sample of $50\,000$ paths. All reported training runs use the same
degree-three basis and the same ridge penalty.

To assess the estimate of Theorem~\ref{th:approx_scheme_error}, we also consider
\[
N\in\{5,10,20,40,80,160\},
\]
or equivalently,
\[
h\in\{0.2,0.1,0.05,0.025,0.0125,0.00625\}.
\]
The discrete processes are reconstructed as left-constant processes on eight
substeps per grid interval. Figures~\ref{fig:example43_convergence},
\ref{fig:example531_convergence}, and \ref{fig:example532_convergence} report
separately
\[
\max_{0\leq t\leq T}\mathbb E|Y_t-\bar Y_t^h|^2
\qquad\text{and}\qquad
\mathbb E\!\int_0^T |Z_t-\bar Z_t^h|^2\,dt.
\]
Each point is averaged over three independent training seeds and evaluated on
$100\,000$ independent paths. For Examples~5.3 (I) and~5.3 (II), both errors
have empirical log--log slopes close to one. In Example~4.3, the $Y$ error has
slope close to one, whereas the $Z$ error decreases faster, with a slope close
to two over most of the retained range. These results are consistent with the
$O(h)$ squared-error estimate of Theorem~\ref{th:approx_scheme_error}; the faster
decay observed for $Z$ in Example~4.3 is specific to this benchmark.

\begin{figure}[ht]
\centering
\includegraphics[width=\textwidth]{Figures/example_4_3_y_z_quantiles.png}
\caption{Example 4.3: exact quantiles (solid lines) and Hermite-3 regression
quantiles (hollow circles) for $Y$ and $Z$.}
\label{fig:example43_quantiles}
\end{figure}

\begin{figure}[ht]
\centering
\includegraphics[width=0.78\textwidth]{Figures/example_4_3_theorem_5_5.png}
\caption{Example 4.3: the two time-reconstruction errors in
Theorem~\ref{th:approx_scheme_error}.}
\label{fig:example43_convergence}
\end{figure}

\begin{figure}[ht]
\centering
\includegraphics[width=\textwidth]{Figures/example_5_3_1_y_z_quantiles.png}
\caption{Example 5.3 (I): exact quantiles (solid lines) and Hermite-3 regression
quantiles (hollow circles) for $Y$ and $Z$.}
\label{fig:example531_quantiles}
\end{figure}

\begin{figure}[ht]
\centering
\includegraphics[width=0.78\textwidth]{Figures/example_5_3_1_theorem_5_5.png}
\caption{Example 5.3 (I): the two time-reconstruction errors in
Theorem~\ref{th:approx_scheme_error}.}
\label{fig:example531_convergence}
\end{figure}

\begin{figure}[ht]
\centering
\includegraphics[width=\textwidth]{Figures/example_5_3_2_y_z_quantiles.png}
\caption{Example 5.3 (II): exact quantiles (solid lines) and Hermite-3 regression
quantiles (hollow circles) for $Y$ and $Z$.}
\label{fig:example532_quantiles}
\end{figure}

\begin{figure}[ht]
\centering
\includegraphics[width=0.78\textwidth]{Figures/example_5_3_2_theorem_5_5.png}
\caption{Example 5.3 (II): the two time-reconstruction errors in
Theorem~\ref{th:approx_scheme_error}.}
\label{fig:example532_convergence}
\end{figure}
